\section{Carbon-Aware Deferral Scheduling}
\label{sec:cads}

\subsection{Problem Formulation}

Given a set of GPU jobs $\mathcal{J}$ and a carbon intensity forecast
$\{c_h\}_{h=0}^{H}$ (gCO\textsubscript{2}/kWh indexed by relative hour),
the deferral scheduling problem is: for each deferrable job $j \in \mathcal{J}$
with original start hour $s_j$ and duration $d_j$ hours, choose a new start
$s_j' \in [s_j,\; s_j + \Delta_{\max}]$ that minimises the mean carbon intensity
over the job's execution window, subject to the hard deadline $\Delta_{\max}$.

The total CO\textsubscript{2} footprint of a schedule is:
\begin{equation}
  \text{CO}_2 = \sum_{j \in \mathcal{J}} g_j \cdot p_j \cdot \eta
    \cdot \frac{1}{1000}\sum_{h=s_j}^{s_j + d_j - 1} c_h
  \label{eq:co2}
\end{equation}
where $g_j$ is the GPU count, $p_j$ the per-GPU TDP in kW, and
$\eta = 1.265$ a combined PUE and cooling overhead.

\subsection{CADS: Baseline Deferral Algorithm}

The Carbon-Aware Deferral Scheduler (CADS) solves the per-job deferral
problem with a vectorised exhaustive search over the candidate window.
For each eligible job we evaluate all $\Delta_{\max}+1$ candidate start hours
in parallel using a cumulative-sum representation of the carbon array,
enabling $\mathcal{O}(1)$ range-mean queries:

\begin{equation}
  \overline{c}(s, d) = \frac{C_{s+d} - C_s}{d}, \quad
  C_h = \sum_{t=0}^{h-1} c_t
  \label{eq:cumsum}
\end{equation}

A job is eligible for deferral if (i) it is flagged \texttt{is\_deferrable}
in the trace metadata and (ii) its original mean carbon intensity exceeds a
configurable threshold $\theta$.  Non-deferrable jobs (online inference,
high-priority) are never shifted.

\subsection{CAPSR: A Domain-Specific Predictor}
\label{sec:capsr}

Standard regression algorithms treat all GPU jobs as draws from a single
distribution.  We propose CAPSR (Carbon-Aware Priority-Stratified Regressor),
a novel domain-specific algorithm built from first principles using only
NumPy and SciPy.  CAPSR introduces three contributions:

\subsubsection{Stratified Sub-Models}

CAPSR partitions the training set into strata defined by the Cartesian product
of \texttt{job\_type\_public} and \texttt{priority\_class}.  For each stratum
$\mathcal{S}_k$ with at least $N_{\min}$ samples, CAPSR fits an independent
weighted Ridge regression:

\begin{equation}
  \hat{\boldsymbol{\beta}}_k
    = \bigl(\mathbf{X}_k^\top \mathbf{W}_k \mathbf{X}_k
             + \alpha \mathbf{I}\bigr)^{-1}
      \mathbf{X}_k^\top \mathbf{W}_k \mathbf{y}_k
  \label{eq:ridge}
\end{equation}

Strata below $N_{\min}$ fall back to the global model.  Because training jobs
and inference jobs differ dramatically in their GPU-hours distributions, this
per-stratum approach outperforms a global model on stratum-level accuracy even
when the overall RMSE is comparable to Random Forest.

\subsubsection{Carbon-Weighted Loss}

When carbon intensity data are available, CAPSR assigns sample weight
\begin{equation}
  w_i = 1 + \frac{c_i}{\bar{c}}, \quad \bar{c} = \frac{1}{N}\sum_i c_i
  \label{eq:weight}
\end{equation}
so that jobs running during high-carbon windows contribute more to the
regression fit, aligning model optimisation with the deferral objective.

\subsubsection{Deferral Confidence Score}

CAPSR exposes a \texttt{predict\_deferral\_confidence} method that returns,
for each job, a per-stratum z-score passed through a sigmoid:
\begin{equation}
  \text{conf}_i = \sigma\!\left(\frac{\hat{y}_i - \mu_k}{\sigma_k}\right),
  \quad \sigma(z) = \frac{1}{1+e^{-z}}
  \label{eq:conf}
\end{equation}
where $\mu_k$ and $\sigma_k$ are the training mean and standard deviation of
stratum $k$.  A confidence near 1 indicates that the job is predicted to
consume above-average GPU-hours \emph{for its type and priority class},
making it a strong carbon-deferral candidate.

This score replaces the static CI threshold of baseline CADS:
jobs with $\text{conf}_i > 0.6$ are prioritised for deferral in
predictive CADS mode (\S\ref{sec:cads}), providing a more principled
ranking than comparing every job to a global threshold.

\subsubsection{Evaluation: Deferral Precision@K}

We introduce Deferral Precision@K as CAPSR's unique evaluation metric:
\begin{equation}
  \text{DP}@K = \frac{|\{i \in \text{top-}K : y_i > \bar{y}\}|}{K}
\end{equation}
where the top-$K$ set is the $K\%$ of test jobs with the highest
\texttt{conf} score.  A random ranker achieves DP@10\% $\approx 0.50$;
CAPSR's stratified confidence should improve this toward 1.0 to the extent
that above-mean GPU-hours jobs are predictable from job metadata alone.

